\begin{abstract}

Charts, figures, and texts derived from data play an important role in policy-making and everyday
decision-making. Fact-checking or understanding such artefacts means understanding how they relate to the data
they purport to be about, which poses a significant challenge due to the complexity involved in tracing back
through the computational pipeline that transformed and manipulated the input. Our paper introduces a new
program analysis framework based on dynamic dependence graphs, which supports the interactive exploration of
the fine-grained relationships between input data and visualisations or other output. We formulate a novel
notion in data provenance, which we call \emph{related inputs}, a relation of mutual relevance or ``cognacy''
between inputs that arises when they contribute to a common feature of the output. We employ Jonsson and
Tarski’s concept of conjugate operators on Boolean algebras to formalise this notion, define a ``related
inputs'' operator as the image function for $R^{-1};\relComp R$, where $R$ is the reachability relation for the graph, , and give a procedure for computing related inputs over a dependence graph.

Additionally, we introduce the “related outputs” operator through the image function for the other composite $R^{−1};\relComp R$. This operator captures a cognacy relation between output features that share common inputs. We demonstrate how to integrate this with prior methods involving Galois connections and execution traces, enhancing efficiency and language agnosticism. This approach requires a single procedure as opposed to separate forward and backward analyses.

Our graphical methodology is not only more rapid in addressing similar problems but is also language-agnostic, broadening its applicability. We showcase various examples where visual outputs are automatically enhanced with interactions supporting “related inputs” queries. Furthermore, we explore the use of projection operators to isolate contributions from specific inputs or demands from particular outputs. This offers a more contextualized and informative query process.

Overall, our framework represents a significant advancement in understanding the intricate relationship between data and its derived visual or textual representations, providing a powerful tool for data analysis and decision-making.\end{abstract}
